% Part: first-order-logic
% Chapter: natural-deduction
% Section: derivations

\documentclass[../../../include/open-logic-section]{subfiles}

\begin{document}

\iftag{FOL}
      {\olfileid{fol}{ntd}{der}}
      {\olfileid{pl}{ntd}{der}}

\olsection{\usetoken{P}{derivation}}

\begin{explain}
We hebben uiteengezet wat een aanname is en de afleidingsregels
gegeven. !!^{derivation}s in natuurlijke deductie worden daaruit
inductief opgebouwd: elke !!{derivation} is ofwel een op zichzelf
staande aanname, ofwel samengesteld uit één, twee of drie
!!{derivation}s gevolgd door een correcte afleidingsstap.
\end{explain}

\begin{defn}[!!^{derivation}]
\Article{derivation} \emph{!!{derivation}} van !!a{sentence}~$!A$ uit
aannamen~$\Gamma$ is een eindige boom van !!{sentence}s die aan de
volgende voorwaarden voldoet:
\begin{enumerate}
\item De bovenste !!{sentence}s van de boom behoren tot $\Gamma$ of
  worden door een afleidingsstap in de boom !!{discharged}.
\item De onderste !!{sentence} van de boom is~$!A$.
\item Elke !!{sentence} in de boom, met uitzondering van de onderste
  zin~$!A$, is een premisse van een correcte toepassing van een
  afleidingsregel waarvan de conclusie direct onder die !!{sentence}
  in de boom staat.
\end{enumerate}
We zeggen dan dat $!A$ de \emph{conclusie} van de !!{derivation} is en
$\Gamma$ de !!{undischarged} aannamen ervan zijn.

Als er !!a{derivation} van $!A$ uit~$\Gamma$ bestaat, zeggen we dat
$!A$ uit~$\Gamma$ \emph{!!{derivable}} is, in symbolen: $\Gamma
\Proves !A$. Als er !!a{derivation} van~$!A$ bestaat waarin elke
aanname is !!{discharged}, schrijven we~$\Proves !A$.
\end{defn}

\begin{ex}
Elke op zichzelf staande aanname is !!a{derivation}. Zo is $!A$ op
zichzelf !!a{derivation}, en $!B$ eveneens. We kunnen daaruit een
nieuwe !!{derivation} verkrijgen door bijvoorbeeld de regel
$\Intro{\land}$ toe te passen:
\begin{prooftree}
\AxiomC{$!A$}
\AxiomC{$!B$}
\RightLabel{\Intro{\land}}
\BinaryInfC{$!A \land !B$}
\end{prooftree}
Deze regels zijn algemeen bedoeld: we kunnen de $!A$ en~$!B$ erin door
willekeurige !!{sentence}s vervangen, bijvoorbeeld door $!C$ en~$!D$.
De conclusie wordt dan $!C \land !D$, en dus is
\begin{prooftree}
\AxiomC{$!C$}
\AxiomC{$!D$}
\RightLabel{\Intro{\land}}
\BinaryInfC{$!C \land !D$}
\end{prooftree}
een correcte !!{derivation}. We kunnen de aannamen ook verwisselen,
zodat $!D$ de rol van~$!A$ vervult en $!C$ die van~$!B$. Derhalve is
\begin{prooftree}
\AxiomC{$!D$}
\AxiomC{$!C$}
\RightLabel{\Intro{\land}}
\BinaryInfC{$!D \land !C$}
\end{prooftree}
eveneens een correcte !!{derivation}.

We kunnen nu een andere regel toepassen, bijvoorbeeld $\Intro{\lif}$.
Daarmee kunnen we een implicatie concluderen en elke aanname
!!{discharge} die identiek is aan het antecedent van die implicatie.
Beide volgende bomen zijn bijgevolg correcte !!{derivation}s:
\begin{prooftree}
\AxiomC{$\Discharge{!C}{1}$}
\AxiomC{$!D$}
\RightLabel{\Intro{\land}}
\BinaryInfC{$!C \land !D$}
\DischargeRule{\Intro{\lif}}{1}
\UnaryInfC{$!C \lif (!C \land !D)$}
\DisplayProof\bottomAlignProof
\AxiomC{$!C$}
\AxiomC{$\Discharge{!D}{1}$}
\RightLabel{\Intro{\land}}
\BinaryInfC{$!C \land !D$}
\DischargeRule{\Intro{\lif}}{1}
\UnaryInfC{$!D \lif (!C \land !D)$}
\end{prooftree}
Zij tonen respectievelijk aan dat $!D \Proves !C \lif (!C \land !D)$
en $!C \Proves !D \lif (!C \land !D)$.

Bedenk dat het intrekken van aannamen een mogelijkheid is, geen
verplichting: we hoeven de aannamen niet in te trekken. We kunnen een
regel in het bijzonder zelfs toepassen als de aannamen niet in de
!!{derivation} voorkomen. Het volgende is bijvoorbeeld toegestaan,
hoewel er geen aanname~$!A$ is om te worden !!{discharged}:
\begin{prooftree}
  \AxiomC{$!B$}
  \DischargeRule{\Intro{\lif}}{1}
  \UnaryInfC{$!A \lif !B$}
\end{prooftree}
\end{ex}
\end{document}
